%!TEX Root = ../main.tex
\chapter{Discrete Valuation Rings and Dedekind Domains}
\label{chap:dvr}
In this chapter, we are going to talk about
a special kind of local rings, discrete valuation rings,
and its non-local analogue, Dedekind domain.

Recall that the (Krull)
\vocab{dimension} of a ring is the longest chain
of prime ideals in a ring $\pp_0\subset\pp_1\subset\dots\pp_d$.
Rings of dimension $0$ are Artininan rings.
The next level of complexity is studying domains of dimension $1$.
In this case, every \textit{nonzero} prime ideal is maximal.
The two main examples to keep in mind are
\begin{itemize}
    \item ring of integers $\mathcal O_K$ in a number field;
    \item $\Spec R$ of an algebraic curve.
\end{itemize}

\section{Discrete Valuation Ring}

\subsection{Definition and Basic Properties}
\begin{definition}[Discrete valuation and DVR]
    Let $K$ be a field. A \vocab{discrete valuation}
    on $K$ a surjective function $v : K \to \ZZ\cup\{\infty\}$
    satisfying the following:
    \begin{itemize}
    \item $v(xy) = v(x)+v(y)$.
    \item $v(x+y) \geq \min(v(x), v(y))$.
    \item $v(x)=\infty$ if and only if $x=0$.
    \end{itemize}
    For any discrete valuation, one defines the ring
    $$V = \{x\in K : v(x)\geq 0\}$$
    (it's easy to check that it forms a ring).
    Every ring arises from discrete valuation this way
    is called a \vocab{discrete valuation ring (DVR)}.
\end{definition}
\begin{example}
Here are two examples from both Number Theory and 
Algebraic geometry that looks similar.
\begin{itemize}
    \item For any prime $p$, the ring 
    $$\ZZ_{(p)} = \left\{\tfrac ab \in\QQ : p\nmid b\right\}$$ 
    is a DVR because it comes with the $p$-adic valuation 
    $v_p : \QQ\to \ZZ\cup\{\infty\}$,
    where $v_p(r)$ is the exponent of $p$ in the prime
    factorization of $r$.
    \item For any field $k$, the ring
    $$k[x]_{(x)} = \left\{\tfrac{f(x)}{g(x)}\in k(x)
    : g(0)\neq 0\right\}$$
    is a DVR. The valuation is define by for each 
    $f\in k[x]$, $v(f)$
    is the exponent of $x$ in factorization of $f$,
    and then extend it by $v(f/g) = v(f)-v(g)$.
    \end{itemize}
\end{example}
\begin{proposition}[Properties of DVR]
\label{prop:dvr_properties}
Let $V$ be a DVR with corresponding valuation
$v : K\to \ZZ\cup\{\infty\}$, where $K=\Frac V$.

Fix an element $\pi\in V$ with $v(\pi)=1$,
then every nonzero ideal of $V$ is generated by $(\pi^k)$
for some integer $k$.
\end{proposition}
Consequence of this facts include
\begin{itemize}
    \item $V$ is a PID (hence Noetherian, UFD, and integrally closed).
    \item $V$ is local ring with unique maximal ideal $(\pi)$.
    \item The only prime ideals of $V$ are $(0)$
    or $(\pi)$.
    \item $V$ has Krull dimension $1$.
\end{itemize}
\begin{proof}
    Let $I$ be a nonzero ideal of $V$.
    Let $m = \min\{v(x) : x\in I\}$, and take an element 
    $x$ such that $v(x)=m$. We claim that $I=(\pi^m)$.

    In one direction, we have to show that $\pi^m\in I$.
    Observe that, $v(\pi^m x^{-1}) = mv(\pi) - v(x) = 0$,
    so $\pi^m x^{-1}\in V$, implying that 
    $x\cdot \pi^mx^{-1} = \pi^m\in I$.

    In another direction, given $y\in I$, then by minimality, we have
    $v(y)\geq m$. Thus, $v(y\cdot \pi^{-m})\geq 0$,
    so $y\cdot \pi^{-m}\in V$, so $y\in (\pi^m)$.
\end{proof}
\subsection{Characterization of DVRs}
We have the following big theorem.
\begin{theorem}[Characterization of DVR]
    \label{thm:dvr}
    Let $(R,\mathfrak m)$ be a local Noetherian domain
    with dimension $1$ (i.e., $\mathfrak m$ is the only nonzero
    prime ideal). 
    Let $k=R/\mm$ be the residue field
    and $K=\Frac R$.
    Then, the following are equivalent.
    \begin{enumerate}[label=(\alph*)]
	\item $R$ is a DVR.
	\item $R$ is integrally closed in $K$.
	\item $\mathfrak m$ is principal.
	\item $\dim_k \mathfrak m/\mathfrak m^2 = 1$.
        \item Every nonzero ideal is a power of $\mm$.
	\item There exists $\pi\in R$ such that
	    all ideals are in form $(0)$ or $(\pi^n)$.
    \end{enumerate}
\end{theorem}
\begin{remark}[DVR, integrally closed, and smoothness]
    If $R$ is the ring obtained by localizing $k[x,y]/(f)$
    at the maximal ideal $(x-a,y-b)$,
    then $\dim R = 1$ (by \Cref{thm:dim_transcendence}).
    Moreover, condition (d) $\dim_k \mm / \mm^2 = 1$ is the easiest
    to check whether $R$ is a DVR.
    Roughly speaking $\mm^2$ mods out the quadratic or higher term,
    so $\mm/\mm^2$ only cares about linear approximation
    of $f$ near point $(a,b)$. 
    Working out the details, it turns out that 
    $$R\text{ is a DVR} \iff \text{at least one of }
    \left.\frac{\partial f}{\partial x}\right|_{(a,b)}
    \text{ and } \left.\frac{\partial f}{\partial y}\right|_{(a,b)}
    \text{ is nonzero},$$
    and this is equivalent to that the curve $f(x,y)=0$
    is smooth at $(a,b)$. In particular,
    \begin{center}
    \fbox{DVR looks like dimension $1$ objects 
    at smooth points.}
    \end{center}
    
    Owing to \Cref{prop:integrally_closed_local}, 
    condition (b) gives an easy condition to check
    whether $k[x,y]/(f)$ is a integrally closed
    when $k$ is algebraically closed;
    one simply needs to evaluate partial derivative 
    of $f$ and show that at no point do both
    $\tfrac{\partial f}{\partial x}$ and $\tfrac{\partial f}{\partial y}$ vanish.
\end{remark}
\begin{proof}
    \textbf{(a) $\Rightarrow$ (b).}
    True because $R$ is a PID, hence UFD,
    and integrally closed (\Cref{prop:ufd_integrally_closed}).

    \textbf{(b) $\Rightarrow$ (c).}
    This is the hardest implication.
    Let $a\neq 0\in \mm$.
    We first claim that
    there exists $n$ such that $\mathfrak m^n\subset (a)$.
    To prove this, $\Spec R/(a)$ has only one point,
    so $R/(a)$ is Artinian. 
    Thus, in $R/(a)$, we have from 
    \Cref{thm:artinian_local_ring} that $\overline{\mathfrak m}^n=0$,
    so $\mathfrak m^n\subseteq (a)$%
    \footnote{Note that we are using the easy implication (b) implies (d) in \Cref{thm:artinian_local_ring}.
    Thus, it is possible to avoid using this theorem.}.
    We pick $n$ minimal so that $\mathfrak m^{n-1}\subsetneq (a)$,
    so there exists $b\in\mathfrak m^{n-1}\setminus (a)$.

    Let $x=\frac ab\in K$. Then, since $b\notin (a)$,
    we have $\frac ba\notin R$,
    so it is not integral.
    \begin{itemize}
	\item If $\frac ba\mathfrak m\subseteq\mathfrak m$,
    then by Cayley-Hamilton's theorem (\Cref{thm:cayley_hamilton})
    on module $\mm$ and the map of multiplication by $\frac ba$,
    we get that $\frac ba$ is integral, a contradiction.
	\item If $\frac ba\mathfrak m\subsetneq\mathfrak m$,
    then, we note that $b\frakm \subseteq \mm^n\subseteq (a)$,
    so $b\mm \subseteq (a)$. Dividing by $a$ both sides
    gives $\tfrac ba\mm\subseteq R$.
    Therefore, $\frac ba\mm$ is an ideal,
    but it is not contained in \textit{the} maximal ideal $\mm$,
    so $\frac ba\mathfrak m = R$. Hence, $\mm = (a/b)$.
    \end{itemize}

    \textbf{(c) $\Rightarrow$ (d).}
    Let $\mathfrak m = (\pi)$.
    Then, $\mathfrak m/\mathfrak m^2 = \Span(\overline\pi)$,
    so $\dim_k\mathfrak m /\mathfrak m^2\leq 1$.
    However, $\mathfrak m / \mathfrak m^2\neq 0$ because if
    $\mm=\mm^2$, then by Nakayama (\Cref{cor:nakayama_local_ring}),
    $\mm=(0)$, contradiction.
    Thus, $\dim(\mm / \mm^2) = 1$.

    \textbf{(d) $\Rightarrow$ (c).}
    We have $\mathfrak m/\mathfrak m^2 = k\overline\pi$.
    Pick $\pi\in\mathfrak m$ lifting $\overline\pi$.
    Then, by Nakayama on lifting generators (\Cref{cor:nakayama_generation}),
    we have $\mathfrak m=(\pi)$.

    \textbf{(c) $\Rightarrow$ (e).}
    Suppose $\mm=(\pi)$.
    Now, given $I\subset R$, we claim that we can find $k$ such that $I\subset (\pi^k)$
    but $I\subsetneq (\pi^{k+1})$.
    If such $k$ does not exist, then $I\subseteq (\pi^k)$
    for all $k$, which means that we have an infinite ascending chain of ideals
    $$I\ \subsetneq\ \pi^{-1}I\ \subsetneq\ \pi^{-2}I\ \subsetneq\  \dots,$$
    contradicting the fact that $R$ is Noetherian.
    
    Having found such $k$, we 
    pick $f\in I\setminus (\pi^{k+1})$, so $f=\pi^kg$
    where $g\notin (\pi)$.
    Thus, $g$ is a unit and so $(f)=(\pi^k)\subset I
    \subset (\pi^k)$, forcing equality,
    so $I=(\pi^k)=\mm^k$.

    \textbf{(e) $\Rightarrow$ (f).}
    If $\mm=\mm^2$, then Nakayama (\Cref{cor:nakayama_local_ring})
    gives $\mm=(0)$, a contradiction.
    Thus, there exists $\pi\in\mm$ but $\pi\notin\mm^2$.
    Let $(\pi) = \mm^r$. Since $\pi\notin\mm^2$, it follows that $r=1$.
    Thus, every ideal is of the form $\mm^n = (\pi^n)$.
    
    \textbf{(f) $\Rightarrow$ (a).}
    We define $v : R\to\ZZ\cup\{\infty\}$ as follows:
    for each $a\in R$, by (e), the ideal $(a)$ must be in form
    $(\pi^e)$ for some integer $e\geq 0$.
    We set $e=v(a)$.
    This $v$ clearly satisfies
    \begin{itemize}
        \item $v(a+b) \geq \min(v(a), v(b))$; and
        \item $v(ab) = v(a)v(b)$ for all $a,b\in R$.
    \end{itemize}
    Now, extend $v$ to $K$ by $v(x/y) = v(x)-v(y)$.
    One can see that the above two properties still hold.
    Finally, we have to show that $a\in R$ if and only if $v(a)\geq 0$.
    The $(\Rightarrow)$ direction is clear.
    For the $(\Leftarrow)$ direction, suppose that 
    $a,b\in R$ such that $v(a/b)\geq 0$, or $v(a)\geq v(b)$.
    Then, $(b)\subseteq (a)$, so $b$ divides $a$, so $a/b\in R$.
    Hence, we conclude that $v$ is a discrete valuation
    with valuation ring $R$.
\end{proof}
%
%
%
%
%
%
\section{Dedekind Domain}
\label{sec:dedekind}
\subsection{Definition and Unique Factorization}
The definition of Dedekind Domain is in the proposition.
\begin{proposition}[Dedekind Domain]
    Let $R$ be a domain, Noetherian, and one-dimensional
    (every nonzero prime ideal is maximal).
    Then, the following are equivalent.
    \begin{enumerate}[label=(\alph*)]
    \item $R$ is integrally closed.
    \item Every localization of nonzero prime  $R_\pp$ is a DVR.
    \item Every primary ideal is a power of a prime ideal.
    \end{enumerate}
    We say that $R$ is \vocab{Dedekind Domain} if
    it meets the (equivalent) conditions of the proposition.
\end{proposition}
Since DVR corresponds to objects at smooth points,
we get that 
\begin{center}
\fbox{Dedekind domain looks like smooth dimension $1$ 
objects.}
\end{center}
\begin{proof}
    \textbf{(a) $\Leftrightarrow$ (b).}
    Since being integrally closed is local (\Cref{prop:integral_closure_local}),
    we get that $R$ is integrally closed 
    if and only if $R_\pp$ is integrally closed for 
    all nonzero prime ideal $\pp$.
    By \Cref{thm:dvr} ((b) implies (a)), this is equivalent to
    $R_\pp$ being DVR for all nonzero prime ideal $\pp$.

    \textbf{(b) $\Rightarrow$ (c).}
    The $(0)$-primary ideal is clearly just $(0)$.
    Thus, it suffices to consider only $\pp$-primary
    for all nonzero prime $\pp$.
    If $Q$ is $\pp$-primary, then
    $Q_\pp$ sits inside a DVR $R_\pp$, so $Q_\pp = \pp^kR_\pp$,
    and so by contracting to $R$, we get that $Q=\pp^k$.

    \textbf{(c) $\Rightarrow$ (b).}
    Now, suppose (c) holds: every primary is a prime power.
    Let $\pp$ be any nonzero prime ideal of $R$,
    and consider any ideal $I\subset R_\pp$.
    Note that the only prime ideals of $I$ and $R_\pp$
    are $(0)$ and $\pp$, so $\Ass(R_\pp/I) = \{\pp\}$
    (it cannot be empty),
    so $I$ is $\pp$-primary.

    Thus, the contraction of $I$ to $R$ is $\pp$-primary,
    and so it must be a power of $\pp$.
    This means that $I = \pp^k$ for some $k$,
    and \Cref{thm:dvr} ((e) implies (a)) concludes that 
    $R_\pp$ is a DVR.
\end{proof}
\begin{corollary}[Dedekind domains has unique factorization of ideals]
    In a Dedekind Domain, any ideal
    $I$ can be factored uniquely to
    $$I = \prod_{i=1}^n \pp_i^{a_i},$$
    where $\pp_i$ are maximal ideals and $a_i\geq 0$.
\end{corollary}
\begin{proof}
    By primary decomposition,
    we have $I=\bigcap_{i=1}^n \pp_i^{a_i}$.
    By Chinese Remainder Theorem,
    it suffices to prove that $\pp_i^{a_i}$ and $\pp_j^{a_j}$
    are coprime.
    But this follows from that $\pp_i$ and $\pp_j$ are coprime.
\end{proof}
\begin{exercise}
Prove that if $R$ is a Dedekind domain, then $R$ is a PID 
if and only if it is a UFD.
\end{exercise}
We now give some examples of Dedekind Domain.
\begin{example}
    \begin{itemize}
    \item PID's are Dedekind Domain because the localization at any nonzero prime ideal is PID, hence a DVR.
    \item The next theorem (\Cref{thm:aklb_dedekind})
    will show that the ring of integers $\calO_K$
    of any number field $K$ is a Dedekind domain.
    This is a key example in number theory.
    % \item If $f(x,y)=0$ cuts out a smooth curve 
    % in the affine plane (over a field $k$), then 
    % $k[x,y]/(f)$ is a Dedekind domain.
    % In other words, ring of smooth curves is a Dedekind domain.
    \end{itemize}
\end{example}
%
%
%
%
\subsection{Key Example: Extensions of Dedekind Domain}
We will now discuss a way to construct a Dedekind Domain.
\begin{theorem}[Extension of Dedekind domain is Dedekind domain]
    \label{thm:aklb_dedekind}
    Let $A\subset K=\Frac A$ be a Dedekind Domain.
    Let $K\subset L$ be a finite, \textbf{separable} extension.
    Let $B$ be an integral closure of $A$ inside $L$.
    Then, $B$ is a Dedekind domain.
\end{theorem}
To prove that $B$ is a Dedekind domain, we need to show that
$B$ is
\begin{itemize}
    \item \textbf{integrally closed}: obvious;
    \item \textbf{one-dimension:} follows from \Cref{thm:dim_integral};
    \item \textbf{Noetherian:} that is (surprisingly) the hard part.
\end{itemize}
The hard part is to prove that $B$ is Noetherian.
This will require the fact that $L\supset K$ is separable.

First, let us recall what separable means.
\begin{definition}[Separable extension]
    A field extension $L/K$ is \vocab{separable} if and only if
    for any $\alpha\in L$, the minimal polynomial
    of $\alpha$ over $K$ has distinct roots in $\overline K\supset K$.
\end{definition}
\begin{example}[Non-seperable extension]
Let $K=\mathbb F_p(t)$ and $L=K[x]/(x^p-t)$.
Then, if $\alpha$ is the image of $x$ in $K[x]/(x^p-t)$,
then the minimal polynomial of $\alpha$ is $x^p-t = (x-\alpha)^p$.
\end{example}

We use the following fact from basic field theory:
``$L/K$ is separable
if and only if there exists $[L:K]=n$ distinct embeddings
$\sigma_1,\dots,\sigma_n : L\to\overline K$''.
If we base change a separable extension 
to algebraic closure, we have
$$\overline K \otimes_K L = \overline K\otimes_K \frac{K[x]}{(f)}
= \frac{\overline K[x]}{(f)}
= \frac{\overline K[x]}{\prod_{i=1}^n (x-\alpha_i)}
\simeq \prod_{i=1}^n \overline K.$$
Then, we have the map
\setlength\arraycolsep{2pt}
$$\begin{array}{rcccl}
    L &\to & \overline K\otimes_K L &\to & \prod_{i=1}^n\overline K \\[4pt]
    y &\mapsto & 1\otimes y &\mapsto & (\sigma_1(y),\dots,\sigma_n(y)).
\end{array}$$

Now, let's go back to \Cref{thm:aklb_dedekind}.
Recall that the tricky part is to prove that $B$ is Noetherian.
This is hard because
\begin{remark}
    If $A$ is Noetherian, then the integral closure
    of $A$ inside $\Frac A$ need not be Noetherian.
\end{remark}
However, we have the following result.
\begin{theorem}
    \label{thm:aklb_finite_module}
    Let $A$ be Noetherian and integrally closed
    inside fraction field $A\subset K$.
    Let $L$ be a finite, separable extension of $K$.
    Let $B$ be the integral closure of $A$ in $L$.
    Then, $B$ is finitely-generated $A$-module.
\end{theorem}
This theorem immediately proves \Cref{thm:aklb_dedekind}.
\begin{proof}[Proof of \Cref{thm:aklb_finite_module}]
    To prove this theorem, we need the notion of trace.
    Define the trace map
    \begin{align*}
	\Tr_{L/K} : L &\to K \\
	y &\mapsto \Tr(\underbrace{z\mapsto yz}_{M_y}).
    \end{align*}
    For separable case, one can calculate trace
    using the basis of $\overline K\otimes L = \prod\overline K$.
    In particular, we have $\Tr(M_y) = \Tr(M_{1\otimes y})$.
    However, if we write this in a basis of $\prod \overline K$,
    we have
    $$M_{1\otimes y} = M_{(\sigma_1(y),\dots,\sigma_n(y))}
    = \begin{pmatrix}
	\sigma_1(y) & \dots & 0 \\
	& \ddots & \\
	0 & \dots & \sigma_n(y)
    \end{pmatrix},$$
    and so we have $\Tr_{L/K}(y) = \sigma_1(y) + \dots + \sigma_n(y)$.

    We now have the following observations.
    \begin{enumerate}
    \item If $b\in B$, then $\Tr_{L/K}(b)\in A$.

    \emph{Proof.} We have $\Tr_{L/K}(b) = \sum_{i=1}^n \sigma_i(b)\in K$.
    Since $b$ is integral, we have $\sigma_i(b)$ integral,
    so $\Tr_{L/K}(b)$ is integral over $A$,
    so $\Tr_{L/K}(b)\in A$.
    \item The trace pairing $\langle \bullet,\bullet\rangle
	: L\times L\to K$ by $\langle y_1,y_2\rangle
	= \Tr_{L/K}(y_1y_2)$ is nondegenerate.
	(i.e., if $y\neq 0$, there exists $z\in L$
	for which $\Tr_{L/K}(yz)\neq 0$).

    \emph{Proof.} We need to check this after tensoring
    with $\overline K$, in which get $\prod K$.
    The pairing then is
    $$\langle (z_1,\dots,z_n), (w_1,\dots,w_n)\rangle
    = z_1w_1+\dots+z_nw_n,$$
    and it's easy to see that this is non-degenerate.
    \end{enumerate}
    
    We now prove the theorem.
    Pick a basis $x_1,\dots,x_n$ of $L$ over $K$.
    By scaling $x_i$ by $a\in A$, one can assume $x_i\in B$
    (just let $a$ be the product of all denominators).

    The next step is to define the trace dual.
    Pick $y_1,\dots,y_n\in L$ such that
    $$\Tr(x_iy_j) = \begin{cases}
	1 & i=j \\ 0 & i\neq j.
    \end{cases}$$
    Then, we are almost one. We claim
    $$B\subset\bigoplus_{i=1}^n Ay_i.$$
    Why? If $b\in B$, write $b=\sum c_iy_i$,
    then by taking trace pairing both sides,
    we get that $\Tr_{L/K}(x_i b) = c_i \in A$,
    so we are done.
    That equation is enough to deduce the theorem.
\end{proof}
Let's give some example applications of this theorem.
\begin{example}
\begin{itemize}
    \item The most obvious application is in number theory.
    When $A=\ZZ$, $K=\QQ$, and $L$ is a finite extension of $\QQ$,
    we get that 
    the ring of integers $B=\calO_L$ is a Dedekind domain.
    \item 
    Let $A=\CC[x]$, $K=\CC(x)$, and $L=\CC(x,y)/(y^2-x^3-x)$.
    Then, one can show that $\CC[x,y]/(y^2-x^3-x)$
    is integrally closed (which is true, but require some work),
    then $B=\CC[x,y]/(y^2-x^3-x)$ is a Dedekind domain.

    In general, smooth curves over $\CC$ correspond 
    to integral closure of a finite extension of $\CC(x)$.
    Then, the integral closure gives a Dedekind domain.
\end{itemize}
\end{example}
%
%
%
%
\section{Fractional Ideals}
\subsection{Fractional Ideals and Invertible Ideals}
We will first discuss these notions for a general domain
(not necessarily Dedekind Domain).
\begin{definition}[Fractional ideals]
    Let $A$ be a domain and $K=\Frac A$.
    An $A$-submodule $I\subset K$ is a \vocab{fractional
    ideal} if there exists $x\in A$ such that $xI\subseteq A$.

    Equivalently, a submodule $I\subset K$ is a fractional ideal
    if and only if it is in form $xJ$ for some $x\in K$
    and regular ideal $J\subset A$.
\end{definition}
\begin{definition}[Invertible and principal fractional ideal]
    A fractional ideal $I\subset K$ is 
    \begin{itemize}
    \item \vocab{invertible}
    if and only if there exists a fractional ideal
    $J\subset K$ such that $IJ=A$.
    \item \vocab{principal}
    if and only if $I=(x)$ for some $x\in K$.
    \end{itemize}
\end{definition}
\begin{proposition}
    If $I$ is invertible, then $I$ is finitely-generated.
\end{proposition}
\begin{proof}
    Suppose that $IJ=A$. Then, there exists
    $i_1,\dots,i_n\in I$ and $j_1,\dots,j_n\in J$
    such that $i_1j_1 + \dots + i_nj_n=1$.
    For any $x\in I$, note that $xj_i \in IJ = A$
    and we have $x = i_1(xj_1) + \dots i_n(xj_n)$,
    so $I$ is generated by $i_1,\dots,i_n$.
\end{proof}
The following theorem shows that invertible fractional
ideals are ``locally principal''.
In other words, they are locally free modules of rank $1$.
\begin{theorem}[Characterization of invertible fractional ideals]
    \label{thm:fractional_ideal}
    Let $I$ be a fractional ideal over domain $A$. 
    Suppose that $I$ is finitely-generated.
    Then, the following
    are equivalent.
    \begin{alphenum}
    \item $I$ is invertible.
    \item $(A:I)I = A$,
    where $(A:I) = \{x\in K : xI\subset A\}$.
    \notation{$(A:I)$}{colon fractional ideal}
    \item $I_\pp$ is an invertible fractional $A_\pp$-ideal 
    for all prime ideal $\pp$.
    \item $I_\mm$ is an invertible fractional $A_\mm$-ideal
    for all maximal ideal $\mm$.
    \item $I_\pp$ is a principal ideal for all prime ideal $\pp$.
    \item $I_\mm$ is a principal ideal for all maximal ideal $\mm$.
    \end{alphenum}
\end{theorem}
\begin{proof}
    We have the following clear implications:
    \begin{itemize}
    \item (b) implies (a) because $(A:I)$ is the inverse.
    \item Localizing $IJ=A$ gives $I_\pp J_\pp = A_\pp$,
    so (a) implies (c), and (c) implies (d).
    \item Principal ideals are invertible, so 
    (e) implies (c) and (f) implies (d).
    \end{itemize}
    We first show (a) implies (b).
    (Thus, (a) and (b) are now equivalent.)
    Suppose $IJ=A$.
    Then, for any $j\in J$, we have $jI \subseteq IJ = A$,
    so $j\in (A:I)$. Thus, $J\subset (A:I)$. 
    In the other direction, note that
    $$(A:I) = (A:I)A = (A:I)IJ \subseteq AJ = J,$$
    so $(A:I)\subseteq J$, forcing $J=(A:I)$.
    
    Now, we show that (a), (b), (c), and (d) are equivalent.
    To do this, it suffices to show that (d) implies (a).
    Assume (d) is true.
    First, we leave to the reader to check that
    $(A:I)_\mm = (A_\mm : I_\mm)$ for all maximal ideals $\mm$.
    (This uses $I$ is finitely-generated.)
    Note that $(A:I)I \subset A$, so we have the inclusion map
    $\phi : (A:I)I \to A$. Localizing this map gives
    $\phi_\mm : (A_\mm : I_\mm)I_\mm \to A_\mm$,
    which we know from (b) (applied on $A_\mm$)
    that it is surjective. 
    Thus, by \Cref{cor:injective_surjective_local},
    we get that $\phi$ is an isomorphism,
    so $(A:I)I = A$.

    Finally, we show that (c) implies (e) and (d) implies (f).
    To show both of these, it suffices to show that
    ``if $A$ is a local ring and 
    $I$ is an invertible fractional ideal, then
    $I$ is principal.''

    To prove this claim, let $\mm$ be the (unique) 
    maximal ideal of $A$, and let $IJ=A$.
    Then, there exists $i_1,\dots,i_n\in I$
    and $j_1,\dots,j_n\in J$ such that 
    $i_1j_1 + \dots + i_nj_n = 1$.
    Thus, at least one of the terms, say $i_1j_1$,
    is in $A\setminus \mm$, hence a unit.
    Now, let $I'=i_1$. Then,
    $$(1) \subseteq (i_1)(j_1) \subseteq (i_1)J
    \subseteq IJ = (1),$$
    forcing all equalities. Thus, $(i_1)J = IJ$.
    Multiplying by $I$ both sides gives $(i_1)=I$.
\end{proof}
\begin{example}
    The ideal $I=(x,y)\subset \CC[x,y]$ is not invertible.
    We check this in two ways.
    \begin{itemize}
    \item First, we show that (b) does not hold.
    We claim that $(\CC[x,y] : I) = \CC[x,y]$.
    Indeed, if $f\in \CC(x,y)$ such that $xf\in \CC[x,y]$
    and $yf\in\CC[x,y]$, then we would have $f\in\CC[x,y]$.
    \item We show that (e) does not hold. To do this,
    we localize at maximal ideal $\mm = (x,y)$. 
    Then, $I_\mm = \mm(\CC[x,y]_\mm) =: \overline\mm$, 
    the maximal ideal of $\CC[x,y]_\mm$. 
    Then, $\overline\mm / \overline\mm^2$ is spanned by 
    $x$ and $y$, so it is definitely not generated by one element,
    so $\overline\mm$ is not principal.
    \end{itemize}
\end{example}
Now, we specialize to Dedekind domain.
\begin{theorem}
    \label{thm:frac_ideal_dedekind}
    If $A$ is a Dedekind domain, then
    every fractional ideal is invertible.
\end{theorem}
\begin{proof}
    Let $I$ be a fractional ideal of $A$.
    Then, since $A_\pp$ is a DVR, hence PID
    (\Cref{thm:dvr}), we get that $I_\pp$ is principal
    for all nonzero prime ideal $\pp$.
    Hence, $I$ is invertible (by \Cref{thm:fractional_ideal}
    (d) implies (a)).
\end{proof}

\subsection{Class Group and Elliptic Curves}
\begin{definition}[Class group]
    Let $A$ be a Dedekind domain. Then, the invertible
    fractional ideals forms the group $\calI_A$
    (the inverse exist because of \Cref{thm:frac_ideal_dedekind}).
    Then, the \vocab{class group} is defined as
    \notation{$\calI_A$}{group of invertible fractional ideal}
    \notation{$\Cl A$}{class group}
    $$\Cl A := \frac{\calI_A}{\{\text{principal fractional
    ideals of } A\}}.$$
\end{definition}
\begin{example}
\begin{itemize}
\item Any PID has a trivial class group.
\item When $\mathcal O_K$ is an integral closure of $\mathbb Z$
inside some field extension of $\mathbb Q$,
the key result in algebraic number theory 
is that $\Cl\calO_K$ is finite.
Proving this require understanding the lattice
structure of $\calO_K$, which is beyond the scope
of these notes.

The \vocab{class number} $\#\Cl\calO_K$ measures
how far away $\mathcal O_K$ is from a UFD.
\item If $A$ be an integral closure of $\mathbb C[x]$
in a finite extension of $\CC(x)$
(thus getting a ring corresponding to a smooth curve),
then there exists a number $g$, called the 
\vocab{genus} of the curve $\Spec A$, such that
$$\Cl A \simeq \CC^g/\Lambda,$$
where $\Lambda\simeq\ZZ^{2g}$ is a lattice
embedded in $\CC^{g}$.
The proof of this is far beyond the scope of these notes.
\end{itemize}
\end{example}

Let's sketch how class group works in a specific case
of \vocab{elliptic curves}. 
(We will not fully justify everything).
Consider the ring
$$A = \frac{\CC[x,y]}{(y^2-x^3-ax-b)},$$
where  $x^3+ax+b$ is a cubic polynomial with distinct roots
(i.e., $4a^3 + 27b^2\neq 0$).
This ring corresponds to elliptic curve $y^2=x^3+ax+b$.
Note that $\dim A=1$ by \Cref{thm:dim_transcendence}.
By using the condition $\dim(\mm/\mm^2)=1$
(in \Cref{thm:dvr} (d)), one can check that localization
$A_\mm$ at any maximal ideal is a DVR, so $A$ is a Dedekind domain.

Now, let us investigate class group.
Suppose that $\mm_{(p,q)} = (x-p, y-q)$
is the maximal ideal corresponding to point $(p,q)$.
We observe the following.
\begin{itemize}
    \item Function $x-p$ vanishes at two points 
    $(p,q)$ and $(p,-q)$. Checking valuation reveals that
    $$(x-p) = \mm_{(p,q)} \mm_{(p,-q)}.$$
    \item If points $(p_1, q_1)$, $(p_2,q_2)$,
    and $(p_3,q_3)$ lie on a line $ux+vy+w=0$.
    Then, because elliptic curve is a cubic equation,
    this line must intersect the curve at
    only these three points.
    Then, checking  valuations reveal that
    $$(ux+vy+w) = \mm_{(p_1,q_1)} \mm_{(p_2,q_2)}
    \mm_{(p_3,q_3)}.$$
\end{itemize}
Thus, if we consider the map 
\begin{align*}
    \phi : \{\text{points on elliptic curve}\}
    \sqcup \{O\} &\to \Cl A \\
    P &\mapsto \mm_P \\
    O &\mapsto (1)
\end{align*}
Then, we get the following observations:
\begin{itemize}
    \item If $-P$ is the reflection of $P$ across the 
    $x$-axis, then $\phi(P)+\phi(-P) = 0$.
    \item If point $P,Q,R$ are collinear, then 
    $\phi(P)+\phi(Q)+\phi(R) = 0$.
    \item In particular, $\phi(P)+\phi(Q)$
    is $\phi$ of the reflection of $R$ across the $x$-axis.
\end{itemize}
\begin{center}
\begin{adjustbox}{valign=m}
\begin{asy}
size(0,7cm);
import graph;
real f(real x){
    real t = x^3-4*x+1;
    return (t>0) ? sqrt(t) : 0;
}
//draw((-2.5,0) -- (2.5,0), gray, EndArrow());
//draw((0,-f(2.5)) -- (0,f(2.5)), gray, EndArrow());
path p = graph(f,-2.1149, 0.25410);
path q = graph(f, 1.8608,3.2);
draw(reverse(p) .. reflect((0,0), (1,0)) * p, linewidth(0.8));
draw(reverse(q) .. reflect((0,0), (1,0)) * q, linewidth(0.8));

pair P = (-1,2);
pair P1 = (-1,-2);
draw((-1,3)--(-1,-3), red+linewidth(0.8));
dot("$P$",P,NE);
dot("$-P$",P1,SE);
label("$\phi(P)+\phi(-P)=0$", (0,-f(3.2)),S);
label("$\mathfrak m_P \mathfrak m_{-P} = (x-p)$", 
(0,-f(3.2)-1),S);
\end{asy}
\end{adjustbox}
\hspace{1.5cm}
\begin{adjustbox}{valign=m}
\begin{asy}
size(0,7cm);
import graph;
real f(real x){
    real t = x^3-4*x+1;
    return (t>0) ? sqrt(t) : 0;
}
//draw((-2.5,0) -- (2.5,0), gray, EndArrow());
//draw((0,-f(2.5)) -- (0,f(2.5)), gray, EndArrow());
path p = graph(f,-2.1149, 0.25410);
path q = graph(f, 1.8608,3.2);
draw(reverse(p) .. reflect((0,0), (1,0)) * p, linewidth(0.8));
draw(reverse(q) .. reflect((0,0), (1,0)) * q, linewidth(0.8));

pair P = (-1,2);
pair Q = (3,4);
pair R = (-7/4, 13/8);
draw(R+0.2*(R-Q)--Q+0.25*(Q-R), red+linewidth(0.8));
label("$ux+vy+w=0$", Q+0.25*(Q-R), N, red);
dot("$P$",P,N);
dot("$Q$",Q,NW);
dot("$R$",R,NW);
label("$\phi(P)+\phi(Q)+\phi(R)=0$", (0,-f(3.2)),S);
label("$\mathfrak m_P \mathfrak m_Q \mathfrak m_R = (ux+vy+wz)$", 
(0,-f(3.2)-1),S);
\end{asy}
\end{adjustbox}
\end{center}
Since $\Cl A$ is generated by maximal ideals,
we have shown that $\phi$ is surjective already.
It turns out that \textbf{$\phi$ is injective}.
This fact is equivalent to $\mm_P$ is not principal for all point $P$,
and proving that requires a nontrivial algebraic geometry
that we will not take time to develop.
(The rough idea is that if $\mm_P$ is generated by $(f)$,
then $f$ vanishes at only point $P$, and we can show that
it induces an isomorphism from $\CC[z]$ to $A$).

The important consequence of this is that
by simply writing $\phi(P)$ as $P$,
we get the group structure on points on elliptic curve
(plus $O$). In particular, if $P,Q,R$ are collinear, then
$P+Q$ is the reflection of $R$ across the $x$-axis.
This is the \vocab{elliptic curve group law},
and notice that we have immediately shown associativity 
of this addition process (which would require checking 
identities with hundreds of terms if bashing it out directly).

This also gives hint on how to generalize to other curves.
The idea is that by considering the set of point 
where a polynomial in variables $x$ and $y$ intersect the curve,
one can show that each element in $\Cl A$
is represented by the product of at most $g$ maximal ideals.
This number $g$ is called the \vocab{genus} of a curve.
Thus, instead of adding points (like elliptic curve),
we can define a group law on $g$-tuple of points instead.
This group law is known as the \vocab{Jacobian} of a curve.
%
%
\section{Problems}
\begin{problem}
Let $A$ be the localization of $\mathbb C[x,y]/(y^3-x^5-x)$ at maximal ideal $(x,y)$.
\begin{alphenum}
\item Prove that $A$ is a DVR.
\item What is the valuation of $x$ in this DVR?
\end{alphenum}
\end{problem}
\begin{problem}
Let $\mathbb Z_{(2)} = \{\tfrac ab : a,b\in\mathbb Z, 2\nmid b\}\subseteq \mathbb Q$. Let $R=\mathbb Z_{(2)}[x]/(x^5+2x+2)$, which is an integral extension of $\mathbb Z_{(2)}$.
\begin{alphenum}
\item Prove that $R$ is a DVR. 
(In particular, you have to show that $R$ is a local ring first.)
\item What is the valuation of $2$ in $R$?
\end{alphenum}
\end{problem}
\begin{problem}
\label{prob:hyperelliptic_dedekind}
Let $f\in\CC[x]$ be a polynomial.
Prove that $\CC[x,y]/(y^2-f(x))$ is a Dedekind domain 
if and only if all roots of $f$ are distinct.
\end{problem}
\begin{problem}
\label{prob:dedekind_ideal}
Let $A$ be a Dedekind domain and $I$ and $J$ be fractional ideals.
Prove that 
\begin{alphenum}
\item $I\oplus J \simeq A\oplus IJ$
\item $I\otimes_A J \simeq IJ$.
\end{alphenum}
This generalizes \Cref{prob:Zsqrt-5}. 
(Hint: localize. See \Cref{sec:locally_free}.)
\end{problem}
\begin{problem}
Let $A$ be a Dedekind domain and $I$ be a nonzero ideal.
\begin{alphenum}
\item Prove that every ideal of $A/I$ is principal.
\item Prove that $I$ is generated by two elements.
\end{alphenum}
\end{problem}
\begin{problem}[2-torsion points of elliptic curves]
Let $\alpha_1,\alpha_2,\alpha_3$ be distinct complex numbers, and let 
$$A=\frac{\mathbb C[x,y]}{(y^2-(x-\alpha_1)(x-\alpha_2)(x-\alpha_3))},$$ 
which is a Dedekind domain by \Cref{prob:hyperelliptic_dedekind}. Prove that the maximal ideal $\mathfrak m_j = (y,x-\alpha_j)$ for $j=1,2,3$ all have order two in the class group $\operatorname{Cl} A$.

(Hint: to show that $\mathfrak m_j$ is not principal, begin by assuming $\frakm_j = (f(x,y))$ and notice that $f(x,-y)$
is also a generator of $\frakm_j$.)
\end{problem}
\begin{problem}[Modules over Dedekind Domain]
In this problem, we will generalize PID structure theorem
(\Cref{thm:pid_structure}) 
to arbitrary an Dedekind domain.
You will find \Cref{prob:projective} and \Cref{sec:locally_free}
helpful.

Let $A$ be a Dedekind Domain and $M$ be an $A$-module.
An element $m\in M$ is a \vocab[torsion element]{torsion} if 
there exists a nonzero element $a\in A$ such that $am=0$.
$M$ is said to be \vocab[torsion module]{torsion}
if every element is a torsion
and $M$ is said to be \vocab{torsion-free}
if the only torsion element of $M$ is $0$.
\begin{alphenum}
\item Prove that any finitely-generated 
torsion $A$-module is a direct sum 
of modules of the form $A/(\pp^{e})$,
where $\pp$ is a prime of $A$ and $e\geq 1$.

(Hint: reduce to the case where $A$ is a DVR by showing that 
$M = \bigoplus_\pp M[\pp^\infty]$,
where $M[\pp^\infty] = \{m\in M : \text{there exists }n 
\text{ and }a\in \pp^n\text{ such that } am=0\}$.)
\item Prove that any finitely-generated torsion-free $A$-module 
is isomorphic to a direct sum of ideals.
Thus, \Cref{prob:dedekind_ideal} concludes that 
every torsion-free $A$-module must be of the form 
$A^{\oplus (r-1)} \oplus I$ for some $r\geq 1$
and fractional ideal $I$.

(Hint: let $K=\Frac A$ and view $M$ as a submodule 
of a $K$-vector space $M\otimes_A K$.)
\item Let $M_{\text{tors}}$ be the submodule of 
torsion elements of $M$.
Prove that $M \simeq M_{\text{tors}} \oplus 
M/M_{\text{tors}}$.
Conclude that any finitely-generated $A$-module 
must be of the form 
$$A/\pp_1^{e_1} \oplus \dots \oplus A/\pp_k^{e_k} 
\oplus A^{\oplus (r-1)} \oplus I,$$
where $\pp_i$ is a prime of $A$, $e_i\geq 0$,
$r\geq 0$, and $I$ is an ideal of $A$.
\end{alphenum}
\end{problem}
\begin{problem}
What about infinitely-generated modules over 
Dedekind domains?
Let $A$ be a Dedekind domain.
\begin{alphenum}
\item (Eilenberg Swindle) Let $I$ be a fractional ideal.
By simplifying 
$$I\oplus I^{-1} \oplus I \oplus I^{-1} \oplus \dots$$
in two ways, show that $A^{\oplus \aleph_0}\oplus I 
\simeq A^{\oplus \aleph_0}$.
\item Prove that for any infinite set $T$ and 
all collections of fractional ideals 
$(I_i)_{i\in T}$ and $(J_i)_{i\in T}$, 
the direct sums $\bigoplus_{i\in T} I_i$ and 
$\bigoplus_{i\in T} J_i$ are isomorphic.
\item Prove that any projective $A$-module 
that is not finitely-generated must be free.
\end{alphenum}
\end{problem}